% Part IV-C — Obligations, Obstructions, and the Construction Loop (Slides 101–120)
\providecommand{\jugmark}[1]{\textbf{\color{jugeoBlue}#1}}

% ---------------------------------------------------------------------------
% Slide 101: What Is a Proof Obligation?
% ---------------------------------------------------------------------------
\begin{frame}
\frametitle{What Is a Proof Obligation?}
\begin{block}{Definition}
  A \jugmark{proof obligation} is something the system needs you (or a solver)
  to prove before it will trust a piece of code.
\end{block}
\vspace{8pt}
\begin{exampleblock}{Example}
  \texttt{def withdraw(balance, amount):}\\
  \quad\texttt{assert amount >= 0}\\[4pt]
  Obligation: \emph{``\texttt{withdraw} always receives a non-negative \texttt{amount}.''}
\end{exampleblock}
\vspace{8pt}
\begin{center}
\begin{tikzpicture}
  \node[fill=jugeoLight, rounded corners, inner sep=8pt, text width=10cm, align=center]
    {\itshape The system \jugmark{generates} obligations; backends \jugmark{discharge} them.};
\end{tikzpicture}
\end{center}
\end{frame}

% ---------------------------------------------------------------------------
% Slide 102: Obligation Lifecycle
% ---------------------------------------------------------------------------
\begin{frame}
\frametitle{Obligation Lifecycle}
\begin{center}
\begin{tikzpicture}[
    state/.style={rectangle, rounded corners=4pt, draw=jugeoBlue, fill=jugeoLight,
                  minimum width=26mm, minimum height=9mm, font=\small\bfseries,
                  line width=1pt, align=center},
    fail/.style={rectangle, rounded corners=4pt, draw=jugeoRed, fill=jugeoRed!10,
                 minimum width=26mm, minimum height=9mm, font=\small\bfseries,
                 line width=1pt, align=center},
    warn/.style={rectangle, rounded corners=4pt, draw=jugeoOrange, fill=jugeoOrange!10,
                 minimum width=26mm, minimum height=9mm, font=\small\bfseries,
                 line width=1pt, align=center},
    arr/.style={-{Stealth[length=5pt]}, line width=1pt},
  ]
  \node[state] (pend)    at (0,0)     {PENDING};
  \node[state] (assign)  at (3.2,0)   {ASSIGNED};
  \node[state] (prog)    at (6.4,0)   {IN\_PROGRESS};
  \node[state] (done)    at (10,0)    {DISCHARGED};
  \node[fail]  (failed)  at (6.4,-1.8){FAILED};
  \node[warn]  (expired) at (3.2,-1.8){EXPIRED};

  \draw[arr, jugeoBlue]   (pend)   -- (assign);
  \draw[arr, jugeoBlue]   (assign) -- (prog);
  \draw[arr, jugeoGreen]  (prog)   -- (done);
  \draw[arr, jugeoRed]    (prog)   -- (failed)
    node[midway, right, font=\scriptsize\color{jugeoRed}] {obstruction};
  \draw[arr, jugeoOrange] (prog)   -- (expired)
    node[midway, above, font=\scriptsize\color{jugeoOrange}] {timeout};
  \draw[arr, jugeoOrange, dashed] (expired) -- (pend)
    node[midway, left, font=\scriptsize\color{jugeoOrange}] {re-schedule};
\end{tikzpicture}
\end{center}
\vspace{6pt}
\begin{itemize}
  \item Happy path: PENDING $\to$ ASSIGNED $\to$ IN\_PROGRESS $\to$ \textcolor{jugeoGreen}{DISCHARGED}
  \item Failure records an \textcolor{jugeoRed}{obstruction}; expiry allows \textcolor{jugeoOrange}{retry}
\end{itemize}
\end{frame}

% ---------------------------------------------------------------------------
% Slide 103: Five Discharge Backends
% ---------------------------------------------------------------------------
\begin{frame}
\frametitle{Five Discharge Backends}
\begin{center}
\small
\begin{tabular}{@{}clll@{}}
  \toprule
  \textbf{\#} & \textbf{Backend} & \textbf{Method} & \textbf{Trust Level}\\
  \midrule
  1 & \jugmark{Structural}  & Vacuous / guaranteed  & \textsc{verified\_proof}\\
  2 & \jugmark{Tautology}   & Pattern matching      & \textsc{verified\_proof}\\
  3 & \jugmark{Z3}          & Fragment-aware SMT    & \textsc{solver\_discharged}\\
  4 & \jugmark{Runtime}     & Test execution        & \textsc{runtime\_witnessed}\\
  5 & \jugmark{Oracle}      & LLM / human           & \textsc{copilot\_suggested}\\
  \bottomrule
\end{tabular}
\end{center}
\vspace{8pt}
\begin{center}
\begin{tikzpicture}
  \node[fill=jugeoLight, rounded corners, inner sep=8pt, text width=10cm, align=center]
    {\itshape Higher in the table $=$ cheaper \textbf{and} more trusted.};
\end{tikzpicture}
\end{center}
\end{frame}

% ---------------------------------------------------------------------------
% Slide 104: Cascade Discharge
% ---------------------------------------------------------------------------
\begin{frame}
\frametitle{Cascade Discharge}
\begin{center}
\begin{tikzpicture}[
    box/.style={rectangle, rounded corners=3pt, draw=jugeoBlue, fill=jugeoLight,
                minimum width=24mm, minimum height=8mm, font=\small\bfseries,
                line width=0.8pt},
    arr/.style={-{Stealth[length=5pt]}, jugeoBlue, line width=0.8pt},
    fail/.style={-{Stealth[length=4pt]}, jugeoRed!60, line width=0.6pt, dashed},
  ]
  \node[box] (s) at (0,0)    {Structural};
  \node[box] (t) at (2.6,0)  {Tautology};
  \node[box] (z) at (5.2,0)  {Z3};
  \node[box] (r) at (7.8,0)  {Runtime};
  \node[box] (o) at (10.4,0) {Oracle};

  \draw[arr]  (s) -- (t) node[midway, above, font=\scriptsize\color{jugeoRed}] {miss};
  \draw[arr]  (t) -- (z) node[midway, above, font=\scriptsize\color{jugeoRed}] {miss};
  \draw[arr]  (z) -- (r) node[midway, above, font=\scriptsize\color{jugeoRed}] {miss};
  \draw[arr]  (r) -- (o) node[midway, above, font=\scriptsize\color{jugeoRed}] {miss};

  \foreach \n in {s,t,z,r,o}
    \draw[-{Stealth[length=4pt]}, jugeoGreen, line width=0.8pt]
      (\n.south) -- ++(0,-0.7)
      node[below, font=\scriptsize\color{jugeoGreen}] {hit!};
\end{tikzpicture}
\end{center}
\vspace{6pt}
\begin{itemize}
  \item Try the \jugmark{cheapest} backend first; fall through on miss
  \item Each hit assigns the \jugmark{appropriate trust level}
  \item Like a medical triage system: easy cases resolved quickly
\end{itemize}
\end{frame}

% ---------------------------------------------------------------------------
% Slide 105: Cohomology Classification of Discharge
% ---------------------------------------------------------------------------
\begin{frame}
\frametitle{Cohomology Classification of Discharge}
\begin{center}
\begin{tikzpicture}[
    cls/.style={rectangle, rounded corners=4pt, minimum width=32mm,
                minimum height=10mm, font=\small\bfseries, line width=1pt,
                align=center},
  ]
  \node[cls, draw=jugeoGreen, fill=jugeoGreen!10]  (h0)  at (0,0)
    {$H^0$ — Fully discharged};
  \node[cls, draw=jugeoOrange, fill=jugeoOrange!10] (h1)  at (0,-1.4)
    {$H^1$ — Partial gap};
  \node[cls, draw=jugeoRed, fill=jugeoRed!10]       (h2)  at (0,-2.8)
    {$H^2$ — Failed};
  \node[cls, draw=gray, fill=gray!10]               (hinf) at (0,-4.2)
    {$H^{-\infty}$ — Deferred to oracle};
\end{tikzpicture}
\end{center}
\vspace{6pt}
\begin{itemize}
  \item $H^0$: global section exists — \textcolor{jugeoGreen}{obligation fully met}
  \item $H^1$: residual gap — \textcolor{jugeoOrange}{partially discharged}
  \item $H^2$: obstruction class — \textcolor{jugeoRed}{proof impossible on current evidence}
  \item $H^{-\infty}$: punted — \textcolor{gray}{awaiting human or LLM input}
\end{itemize}
\end{frame}

% ---------------------------------------------------------------------------
% Slide 106: Comparison — F* vs JuGeo Obligation Handling
% ---------------------------------------------------------------------------
\begin{frame}
\frametitle{F* vs JuGeo — Obligation Handling}
\begin{columns}[T]
\begin{column}{0.47\textwidth}
\begin{alertblock}{F* Pipeline}
  \begin{enumerate}
    \item VC generated
    \item Sent to Z3
    \item \texttt{sat} / \texttt{unsat} / \texttt{unknown}
  \end{enumerate}
  \vspace{4pt}
  {\small One backend. Binary outcome.\\No trust annotation.}
\end{alertblock}
\end{column}
\begin{column}{0.47\textwidth}
\begin{exampleblock}{JuGeo Pipeline}
  \begin{enumerate}
    \item VC generated
    \item \jugmark{Classified} by type
    \item \jugmark{Routed} to best backend
    \item Multi-backend \jugmark{cascade}
    \item Trust-annotated result
    \item Obstruction if failed
  \end{enumerate}
\end{exampleblock}
\end{column}
\end{columns}
\vspace{6pt}
\begin{center}
\begin{tikzpicture}
  \node[fill=jugeoLight, rounded corners, inner sep=6pt, text width=10cm, align=center,
        font=\small]
    {\itshape JuGeo never says just ``unknown'' — it tells you \jugmark{why} and \jugmark{what to do next}.};
\end{tikzpicture}
\end{center}
\end{frame}

% ---------------------------------------------------------------------------
% Slide 107: What Is an Obstruction?
% ---------------------------------------------------------------------------
\begin{frame}
\frametitle{What Is an Obstruction?}
\begin{block}{Key idea}
  When a proof \textcolor{jugeoRed}{fails}, the failure itself becomes a
  \jugmark{first-class object}.
\end{block}
\vspace{8pt}
\begin{columns}[T]
\begin{column}{0.47\textwidth}
\begin{alertblock}{Traditional}
  \begin{itemize}
    \item Error message printed
    \item Scrolls off screen
    \item No structure, no memory
  \end{itemize}
\end{alertblock}
\end{column}
\begin{column}{0.47\textwidth}
\begin{exampleblock}{JuGeo}
  \begin{itemize}
    \item \jugmark{Obstruction} object created
    \item Stored in the registry
    \item Queryable, persistent, repairable
  \end{itemize}
\end{exampleblock}
\end{column}
\end{columns}
\vspace{8pt}
\begin{center}
\begin{tikzpicture}
  \node[fill=jugeoLight, rounded corners, inner sep=6pt, font=\small]
    {\itshape Not an error message — a structured mathematical record.};
\end{tikzpicture}
\end{center}
\end{frame}

% ---------------------------------------------------------------------------
% Slide 108: Anatomy of an Obstruction
% ---------------------------------------------------------------------------
\begin{frame}[fragile]
\frametitle{Anatomy of an Obstruction}
\begin{center}
\begin{tikzpicture}[
    field/.style={rectangle, rounded corners=2pt, draw=jugeoBlue, fill=jugeoLight,
                  minimum width=50mm, minimum height=7mm, font=\ttfamily\scriptsize,
                  line width=0.6pt, anchor=west},
    lab/.style={font=\scriptsize\bfseries\color{jugeoDark}, anchor=east},
  ]
  \node[lab] at (-0.2, 0)    {obstruction\_id};
  \node[field] at (0, 0)     {``obs-042''};
  \node[lab] at (-0.2,-0.9)  {coordinate};
  \node[field] at (0,-0.9)   {where in the codebase};
  \node[lab] at (-0.2,-1.8)  {proposition};
  \node[field] at (0,-1.8)   {what was supposed to hold};
  \node[lab] at (-0.2,-2.7)  {admissibility};
  \node[field] at (0,-2.7)   {which law was violated};
  \node[lab] at (-0.2,-3.6)  {evidence\_present};
  \node[field] at (0,-3.6)   {True / False};
  \node[lab] at (-0.2,-4.5)  {repair\_frontier};
  \node[field, draw=jugeoGreen, fill=jugeoGreen!10] at (0,-4.5) {concrete fix suggestions};
  \node[lab] at (-0.2,-5.4)  {cohomology\_class};
  \node[field] at (0,-5.4)   {H0 / H1 / H2};
  \node[lab] at (-0.2,-6.3)  {downstream};
  \node[field] at (0,-6.3)   {obligations that depend on this};
\end{tikzpicture}
\end{center}
\end{frame}

% ---------------------------------------------------------------------------
% Slide 109: Obstructions Are Persistent
% ---------------------------------------------------------------------------
\begin{frame}
\frametitle{Obstructions Are Persistent}
\begin{center}
\begin{tikzpicture}[
    ses/.style={rectangle, rounded corners=3pt, draw=jugeoBlue, fill=jugeoLight,
                minimum width=28mm, minimum height=9mm, font=\small\bfseries,
                line width=0.8pt},
    obs/.style={circle, draw=jugeoRed, fill=jugeoRed!15, minimum size=6mm,
                font=\scriptsize\bfseries, line width=0.8pt},
    arr/.style={-{Stealth[length=4pt]}, jugeoBlue!60, line width=0.6pt},
  ]
  \node[ses] (s1) at (0,0)   {Session 1};
  \node[ses] (s2) at (3.6,0) {Session 2};
  \node[ses] (s3) at (7.2,0) {Session 3};

  \node[obs] (o1) at (0,-1.4)   {obs};
  \node[obs] (o2) at (1.2,-1.4) {obs};
  \node[obs] (o3) at (3.6,-1.4) {obs};
  \node[obs] (o4) at (4.8,-1.4) {obs};
  \node[obs] (o5) at (6,-1.4)   {obs};

  \draw[arr] (s1)--(o1); \draw[arr] (s1)--(o2);
  \draw[arr] (o2)--(s2); \draw[arr] (s2)--(o3); \draw[arr] (s2)--(o4);
  \draw[arr] (o4)--(s3); \draw[arr] (o5)--(s3);

  \node[font=\footnotesize\itshape\color{jugeoDark}, text width=8cm, align=center]
    at (3.6,-2.6) {Obstructions survive across sessions and accumulate in a registry.};
\end{tikzpicture}
\end{center}
\vspace{4pt}
\begin{itemize}
  \item Unlike error messages that \textcolor{jugeoRed}{scroll away}
  \item Obstructions \jugmark{persist}, \jugmark{accumulate}, and form a structured history
  \item You can return days later and pick up exactly where you left off
\end{itemize}
\end{frame}

% ---------------------------------------------------------------------------
% Slide 110: Obstructions Are Queryable
% ---------------------------------------------------------------------------
\begin{frame}[fragile]
\frametitle{Obstructions Are Queryable}
\begin{lstlisting}
# All obstructions below SOLVER_DISCHARGED trust
low_trust = registry.query(
    trust_below="SOLVER_DISCHARGED"
)

# All obstructions in the auth module
auth_obs = registry.query(
    coordinate_prefix="myapp.auth"
)

# Obstructions that have repair frontiers
repairable = registry.query(
    has_repair_frontier=True
)
\end{lstlisting}
\vspace{6pt}
\begin{center}
\begin{tikzpicture}
  \node[fill=jugeoLight, rounded corners, inner sep=6pt, font=\small, text width=10cm,
        align=center]
    {\itshape Obstructions are a \jugmark{database}, not a log file.};
\end{tikzpicture}
\end{center}
\end{frame}

% ---------------------------------------------------------------------------
% Slide 111: Obstructions Have Repair Frontiers
% ---------------------------------------------------------------------------
\begin{frame}
\frametitle{Obstructions Have Repair Frontiers}
\begin{block}{What is a repair frontier?}
  Each obstruction carries \jugmark{concrete suggestions} for how to resolve it.
\end{block}
\vspace{6pt}
\begin{exampleblock}{Example Repair Frontier}
  \begin{itemize}
    \item ``Check edge case: empty list''
    \item ``Verify base case \texttt{total=0} satisfies spec''
    \item ``Add precondition \texttt{len(items) > 0}''
  \end{itemize}
\end{exampleblock}
\vspace{6pt}
\begin{columns}[T]
\begin{column}{0.47\textwidth}
\begin{alertblock}{Traditional}
  \texttt{Error: assertion failed}\\[2pt]
  {\small Opaque. Now what?}
\end{alertblock}
\end{column}
\begin{column}{0.47\textwidth}
\begin{exampleblock}{JuGeo}
  Structured repair steps\\[2pt]
  {\small Actionable. Try these first.}
\end{exampleblock}
\end{column}
\end{columns}
\end{frame}

% ---------------------------------------------------------------------------
% Slide 112: Obstruction Example
% ---------------------------------------------------------------------------
\begin{frame}[fragile]
\frametitle{Obstruction Example}
\begin{lstlisting}
Obstruction(
  obstruction_id = "obs-spec-affine-007",
  coordinate     = "spec.affine.program.00",
  proposition    = "spec(solve([], -2, 3, 0), ...) == True",
  admissibility  = "GLUING_LAW",
  evidence       = False,
  repair_frontier = [
      "Check edge case: empty list",
      "Verify base case total=0 satisfies spec",
  ],
  cohomology_class = "H1",
  downstream = ["obl-affine-012", "obl-affine-013"],
)
\end{lstlisting}
\vspace{4pt}
\begin{center}
\begin{tikzpicture}
  \node[fill=jugeoOrange!15, rounded corners, inner sep=6pt, font=\small,
        text width=10cm, align=center]
    {$H^1$: partial gap — evidence is missing but a \jugmark{repair path} exists.};
\end{tikzpicture}
\end{center}
\end{frame}

% ---------------------------------------------------------------------------
% Slide 113: Backpressure — When Obstructions Pile Up
% ---------------------------------------------------------------------------
\begin{frame}
\frametitle{Backpressure — When Obstructions Pile Up}
\begin{center}
\begin{tikzpicture}[
    bp/.style={rectangle, rounded corners=3pt, minimum width=42mm,
               minimum height=8mm, font=\small\bfseries, line width=0.8pt,
               align=center},
  ]
  \node[bp, draw=jugeoRed, fill=jugeoRed!10]       at (0,0)
    {OBLIGATION\_OVERFLOW};
  \node[bp, draw=jugeoOrange, fill=jugeoOrange!10]  at (0,-1.1)
    {INTEGRATION\_LAG};
  \node[bp, draw=jugeoOrange, fill=jugeoOrange!10]  at (0,-2.2)
    {TREATY\_INSTABILITY};
  \node[bp, draw=jugeoRed, fill=jugeoRed!10]        at (0,-3.3)
    {EVIDENCE\_EXHAUSTION};
  \node[bp, draw=jugeoRed!80!black, fill=jugeoRed!20] at (0,-4.4)
    {BUDGET\_CRITICAL};
\end{tikzpicture}
\end{center}
\vspace{4pt}
\begin{itemize}
  \item Too many unresolved obligations? The system notices.
  \item Each family triggers a different \jugmark{response strategy}
  \item No other proof assistant has backpressure management
\end{itemize}
\end{frame}

% ---------------------------------------------------------------------------
% Slide 114: Backpressure Responses
% ---------------------------------------------------------------------------
\begin{frame}
\frametitle{Backpressure Responses}
\begin{center}
\small
\begin{tabular}{@{}lp{7.5cm}@{}}
  \toprule
  \textbf{Response} & \textbf{What Happens}\\
  \midrule
  \jugmark{THROTTLE}   & Slow down: reduce obligation generation rate\\
  \jugmark{PAUSE}      & Stop: halt verification until pressure drops\\
  \jugmark{REDIRECT}   & Reroute: try a different backend or strategy\\
  \jugmark{ESCALATE}   & Ask for help: flag for human / LLM review\\
  \jugmark{SHED\_LOAD} & Drop low-priority obligations to free capacity\\
  \bottomrule
\end{tabular}
\end{center}
\vspace{8pt}
\begin{center}
\begin{tikzpicture}
  \node[fill=jugeoLight, rounded corners, inner sep=8pt, text width=10cm, align=center,
        font=\small]
    {\itshape Think of it as \jugmark{congestion control} for proofs —
     the same idea that keeps the internet from collapsing.};
\end{tikzpicture}
\end{center}
\end{frame}

% ---------------------------------------------------------------------------
% Slide 115: The Construction Loop — Overview
% ---------------------------------------------------------------------------
\begin{frame}
\frametitle{The Construction Loop — Overview}
\begin{block}{JuGeo's core verification engine}
  Four phases that repeat until all obligations are discharged (or deferred).
\end{block}
\vspace{6pt}
\begin{center}
\begin{tikzpicture}[
    phase/.style={rectangle, rounded corners=5pt, draw=jugeoBlue, fill=jugeoLight,
                  minimum width=28mm, minimum height=11mm, font=\bfseries,
                  line width=1.2pt, align=center},
    arr/.style={-{Stealth[length=6pt]}, jugeoBlue, line width=1.2pt},
  ]
  \node[phase] (prop) at (0,0)     {PROPOSE};
  \node[phase] (norm) at (3.6,0)   {NORMALIZE};
  \node[phase] (comp) at (7.2,0)   {COMPARE};
  \node[phase] (sel)  at (10.8,0)  {SELECT};

  \draw[arr] (prop) -- (norm);
  \draw[arr] (norm) -- (comp);
  \draw[arr] (comp) -- (sel);
  \draw[arr, rounded corners=8pt] (sel.south) -- ++(0,-0.9) -| (prop.south);
\end{tikzpicture}
\end{center}
\vspace{6pt}
\begin{itemize}
  \item Analogous to F*'s VC generation + discharge, but \jugmark{generalized}
  \item Multiple candidates compete; the best one wins each round
\end{itemize}
\end{frame}

% ---------------------------------------------------------------------------
% Slide 116: Phase 1 — PROPOSE
% ---------------------------------------------------------------------------
\begin{frame}
\frametitle{Phase 1 — PROPOSE}
\begin{center}
\begin{tikzpicture}[
    src/.style={rectangle, rounded corners=3pt, draw=jugeoBlue, fill=jugeoLight,
                minimum width=24mm, minimum height=9mm, font=\small\bfseries,
                line width=0.8pt},
    cand/.style={rectangle, rounded corners=3pt, draw=jugeoGreen, fill=jugeoGreen!10,
                 minimum width=20mm, minimum height=8mm, font=\small,
                 line width=0.8pt},
    arr/.style={-{Stealth[length=5pt]}, jugeoBlue!70, line width=0.7pt},
  ]
  \node[src] (z3)  at (-3, 1)  {Z3};
  \node[src] (rt)  at (-3, 0)  {Runtime};
  \node[src] (cp)  at (-3,-1)  {Copilot};
  \node[src] (hu)  at (-3,-2)  {Human};

  \node[cand] (c1) at (2, 0.8) {candidate A};
  \node[cand] (c2) at (2,-0.2) {candidate B};
  \node[cand] (c3) at (2,-1.2) {candidate C};
  \node[cand] (c4) at (2,-2.2) {candidate D};

  \draw[arr] (z3.east)  -- (c1.west);
  \draw[arr] (rt.east)  -- (c2.west);
  \draw[arr] (cp.east)  -- (c3.west);
  \draw[arr] (hu.east)  -- (c4.west);

  \node[font=\footnotesize\itshape\color{jugeoDark}] at (2, 1.6)
    {Candidate Pool};
\end{tikzpicture}
\end{center}
\vspace{6pt}
\begin{itemize}
  \item Solicit candidate solutions from \jugmark{all channels} simultaneously
  \item Multiple candidates \jugmark{compete} — no single point of failure
  \item Even partial solutions are welcome (they carry partial evidence)
\end{itemize}
\end{frame}

% ---------------------------------------------------------------------------
% Slide 117: Phase 2 — NORMALIZE
% ---------------------------------------------------------------------------
\begin{frame}
\frametitle{Phase 2 — NORMALIZE}
\begin{columns}[T]
\begin{column}{0.47\textwidth}
\begin{alertblock}{Before normalization}
  \ttfamily\scriptsize
  candidate A: \textbackslash{}x -> x + 1\\
  candidate B: lambda y: y + 1\\
  candidate C: def f(z): return z+1
\end{alertblock}
\end{column}
\begin{column}{0.47\textwidth}
\begin{exampleblock}{After normalization}
  \ttfamily\scriptsize
  candidate A: $\lambda v_0.\; v_0 + 1$\\
  candidate B: $\lambda v_0.\; v_0 + 1$\\
  candidate C: $\lambda v_0.\; v_0 + 1$
\end{exampleblock}
\end{column}
\end{columns}
\vspace{10pt}
\begin{block}{What normalization does}
  \begin{itemize}
    \item \jugmark{Alpha-equivalence}: rename variables to canonical forms
    \item \jugmark{Formatting}: strip whitespace and stylistic differences
    \item \jugmark{Desugar}: expand syntactic sugar to core forms
  \end{itemize}
\end{block}
\vspace{4pt}
\begin{center}
  {\small\itshape So candidates can be \jugmark{fairly compared} on substance, not style.}
\end{center}
\end{frame}

% ---------------------------------------------------------------------------
% Slide 118: Phase 3 — COMPARE
% ---------------------------------------------------------------------------
\begin{frame}
\frametitle{Phase 3 — COMPARE}
\begin{center}
\begin{tikzpicture}[
    crit/.style={rectangle, rounded corners=3pt, draw=jugeoBlue, fill=jugeoLight,
                 minimum width=30mm, minimum height=8mm, font=\small\bfseries,
                 line width=0.8pt},
    arr/.style={-{Stealth[length=4pt]}, jugeoBlue!60, line width=0.6pt},
  ]
  \node[crit] (t) at (0, 0)    {Trust Level};
  \node[crit] (r) at (4, 0)    {Residual Count};
  \node[crit] (e) at (8, 0)    {Evidence Strength};

  \node[rectangle, rounded corners=5pt, draw=jugeoGreen, fill=jugeoGreen!10,
        minimum width=30mm, minimum height=9mm, font=\bfseries,
        line width=1pt] (pareto) at (4,-1.8) {Pareto Ranking};

  \draw[arr] (t.south) -- (pareto.north west);
  \draw[arr] (r.south) -- (pareto.north);
  \draw[arr] (e.south) -- (pareto.north east);
\end{tikzpicture}
\end{center}
\vspace{8pt}
\begin{itemize}
  \item \jugmark{Trust level}: \textsc{verified\_proof} $>$ \textsc{solver\_discharged}
        $>$ \textsc{runtime\_witnessed} $>$ \textsc{copilot\_suggested}
  \item \jugmark{Residual count}: fewer unresolved obligations $=$ better
  \item \jugmark{Evidence strength}: how much of the spec is covered?
  \item Multi-criteria $\Rightarrow$ \jugmark{Pareto ranking}, not a single score
\end{itemize}
\end{frame}

% ---------------------------------------------------------------------------
% Slide 119: Phase 4 — SELECT
% ---------------------------------------------------------------------------
\begin{frame}
\frametitle{Phase 4 — SELECT}
\begin{center}
\begin{tikzpicture}[
    cand/.style={rectangle, rounded corners=3pt, minimum width=22mm,
                 minimum height=8mm, font=\small, line width=0.8pt},
    arr/.style={-{Stealth[length=5pt]}, line width=0.8pt},
  ]
  \node[cand, draw=gray, fill=gray!10]              (a) at (0, 1)  {candidate A};
  \node[cand, draw=jugeoGreen, fill=jugeoGreen!15,
        font=\small\bfseries]                        (b) at (0, 0)  {candidate B \checkmark};
  \node[cand, draw=gray, fill=gray!10]              (c) at (0,-1)  {candidate C};

  \node[rectangle, rounded corners=5pt, draw=jugeoGreen, fill=jugeoGreen!10,
        minimum width=34mm, minimum height=10mm, font=\bfseries,
        line width=1.2pt] (win) at (5.5, 0) {Verified Solution};

  \draw[arr, jugeoGreen] (b.east) -- (win.west);

  \node[rectangle, rounded corners=3pt, draw=jugeoBlue, fill=jugeoLight,
        font=\scriptsize, text width=34mm, align=center]
    (just) at (5.5,-1.6)
    {Justification recorded:\\trust, residual, evidence};

  \draw[-{Stealth[length=4pt]}, jugeoBlue!50, line width=0.5pt, dashed]
    (win.south) -- (just.north);
\end{tikzpicture}
\end{center}
\vspace{6pt}
\begin{itemize}
  \item Pick the Pareto-optimal candidate
  \item \jugmark{Explicit justification} recorded — auditable choice
  \item Winner becomes the verified solution for this coordinate
  \item Losers are kept as fallback evidence
\end{itemize}
\end{frame}

% ---------------------------------------------------------------------------
% Slide 120: Compression Records
% ---------------------------------------------------------------------------
\begin{frame}
\frametitle{Compression Records}
\begin{block}{Each construction round produces a delta record}
  A compact summary of \jugmark{what changed} this round.
\end{block}
\vspace{6pt}
\begin{center}
\begin{tikzpicture}[
    row/.style={rectangle, rounded corners=2pt, draw=jugeoBlue, fill=jugeoLight,
                minimum width=48mm, minimum height=7mm, font=\small,
                line width=0.6pt, anchor=west},
    lab/.style={font=\scriptsize\bfseries\color{jugeoDark}, anchor=east},
  ]
  \node[lab] at (-0.2, 0)    {section\_deltas};
  \node[row] at (0, 0)       {Which sections were updated};
  \node[lab] at (-0.2,-0.9)  {obligation\_deltas};
  \node[row] at (0,-0.9)     {New, discharged, or expired obligations};
  \node[lab] at (-0.2,-1.8)  {evidence\_deltas};
  \node[row] at (0,-1.8)     {New evidence collected this round};
  \node[lab] at (-0.2,-2.7)  {obstruction\_deltas};
  \node[row] at (0,-2.7)     {Obstructions added or resolved};
  \node[lab] at (-0.2,-3.6)  {certificate\_updates};
  \node[row] at (0,-3.6)     {Certificate version incremented};
\end{tikzpicture}
\end{center}
\vspace{6pt}
\begin{center}
\begin{tikzpicture}
  \node[fill=jugeoLight, rounded corners, inner sep=6pt, font=\small,
        text width=10cm, align=center]
    {\itshape Full traceability: you can replay the entire verification history
     from compression records alone.};
\end{tikzpicture}
\end{center}
\end{frame}
